\documentclass[spanish,letterpaper,oneside]{article}

\def\documenttitle {Tercer Momento: Propuesta de Innovación Matemática «Calibrar y Comparar»}
\def\documentsubtitle {Fundamentos para la enseñanza y el aprendizaje I · Subproducto 6C (Bloque III)}
\def\documentsubject {Actividad 6 · Tercer Momento}
\def\documentauthor {Mtro. Martín Jonathan de la Cruz Muñoz}
\def\coursename {Fundamentos para la enseñanza y el aprendizaje I}
\def\coursecode {MAEM-SD26}

\def\universityname {Instituto de Investigación, Innovación y Estudios de Posgrado para la Educación del Estado de Nuevo León}
\def\universityfaculty {Maestría en Aprendizaje y Enseñanza de las Matemáticas}
\def\universitydepartment {Fundamentos para la enseñanza y el aprendizaje I}
\def\universitydepartmentimage {img/iiiepe/marca-iiiepe.png}
\def\universitydepartmentimagecfg {height=1.75cm}
\def\universitylocation {Monterrey, Nuevo León}
\def\studentfullname {Mtro. Martín Jonathan de la Cruz Muñoz}
\def\studentid {00874}
\def\studentemail {martin.delacruz@campus.iiiepe.edu.mx}
\def\programperiod {2do periodo, septiembre--diciembre de 2026 (SD26)}
\def\programmodality {Modalidad presencial}
\def\programshort {MAEM}

\def\authortable {
  \begin{tabular}{ll}
    \textbf{Alumno:} & \studentfullname \\
    Matrícula: & \studentid \\
    Correo institucional: & \studentemail \\
    Unidad: & \coursename \\
    Programa: & \programshort \\
    Periodo: & \programperiod \\
    & \\
    \multicolumn{2}{l}{Fecha de realización: 19 de noviembre de 2026} \\
    \multicolumn{2}{l}{\universitylocation}
  \end{tabular}
}

\input{template}

\usepackage{helvet}
\renewcommand{\familydefault}{\sfdefault}
\usepackage{array}
\usepackage{longtable}
\usepackage{pdflscape}
\usepackage{tikz}
\usepackage{pgfplots}
\pgfplotsset{compat=1.18}
\usepackage[most]{tcolorbox}
\usetikzlibrary{arrows.meta,positioning,calc,fit,backgrounds,shapes.geometric}

\definecolor{iiiepeNavy}{HTML}{163B5C}
\definecolor{iiiepeBlue}{HTML}{1F6F9F}
\definecolor{iiiepeAqua}{HTML}{20B7A8}
\definecolor{iiiepeOrange}{HTML}{F28C28}
\definecolor{iiiepePaper}{HTML}{F6F8F7}
\definecolor{iiiepeInk}{HTML}{1E2A32}
\definecolor{iiiepePale}{HTML}{EEF5F8}
\definecolor{iiiepeMint}{HTML}{E7F6F4}
\definecolor{iiiepeWarm}{HTML}{FFF3E6}
\def\portraittitlecolor {iiiepeNavy}
\setcitestyle{authoryear,open={(},close={)}}
\makeatletter
\renewcommand{\@biblabel}[1]{}
\makeatother
\setlength{\bibhang}{1.27cm}
\setlength{\bibsep}{6pt}
\AtBeginDocument{\renewcommand{\refname}{Referencias}}

\newcommand{\insertcoverwatermark}{%
  \AddToShipoutPictureBG*{%
    \begin{tikzpicture}[remember picture,overlay]
      \node[opacity=0.04,inner sep=0pt] at (current page.center) {%
        \includegraphics[width=0.65\paperwidth]{img/iiiepe/marca-iiiepe.png}%
      };
    \end{tikzpicture}%
  }%
}

\setlength{\parskip}{5pt}\setlength{\parindent}{0pt}
\begin{document}\insertcoverwatermark\templatePortrait\templatePagecfg\templateFinalcfg\fancyfoot{}\fancyfoot[L]{\scriptsize IIIEPE · MAEM · SD26}\fancyfoot[C]{\scriptsize\thepage}\fancyfoot[R]{\scriptsize Subproducto 6C · Tercer Momento}\setstretch{1.3}\raggedbottom
\section{Tercer momento: diseño de propuesta de innovación pedagógica}
Con sustento en el análisis teórico del aprendizaje profundo \citep{quiroga2022aprendizaje} y el diagnóstico propio de enseñanza, presento como tercer momento de articulación entre teoría y práctica una propuesta de innovación didáctica \textbf{por realizar} denominada \textit{«Calibrar y comparar para decidir»}.

Mi postura es que transformar el aula técnica e ingenieril requiere una intervención viable, honesta —diferenciando explícitamente un diseño proyectado y datos didácticos simulados de resultados reales futuros— y estructurada alrededor de las dimensiones cognitiva, interpersonal e intrapersonal del estudiantado \citep{fullan2014rich,rivas2017innovaciones}.

\section{Los siete elementos de la propuesta de innovación (por realizar)}
\begin{tcolorbox}[colback=iiiepePale,colframe=iiiepeNavy,title=\textbf{Ficha técnica integral · «Calibrar y comparar para decidir» (por realizar)},breakable]
\begin{enumerate}
  \item \textbf{Contenido matemático:} Modelación lineal ($y = mx + b$), contraste con proporcionalidad directa ($y = kx$), tasa de cambio, ordenada al origen, punto de equilibrio e intervalo físico de validez en cuarto semestre de bachillerato técnico en mecatrónica/mantenimiento e ingeniería inicial \citep{sep2023pensamiento,conalep2025mecatronica}.
  \item \textbf{Problemática atendida:} Desconexión entre manipulación simbólica automática e interpretación física, dimensional y argumentativa de modelos lineales.
  \item \textbf{Propósito de aprendizaje:} Lograr que el estudiantado formule, compare y defienda modelos lineales desde tablas simuladas o experimentales, integrando representaciones numérica, gráfica y algebraica.
  \item \textbf{Situación o desafío matemático:} Dos sensores térmicos ($S_1$ calibrado $y=x$ y $S_2$ simulado $y=1.05x+0.2$) y dos esquemas tarifarios de servicio técnico ($C_A(h)=250+80h$ y $C_B(h)=100+110h$) generan decisiones distintas según el rango de operación. Los equipos predicen sin fórmula inicial, determinan el punto de cruce ($h=5\text{ h}$), discuten el error de suponer proporcionalidad directa y redactan un dictamen justificado con datos didácticos simulados.
  \item \textbf{Estrategias de enseñanza:} Asociación docente--estudiante en tres sesiones: predicción individual inicial, modelación colaborativa en tríos con roles rotativos (\textit{modelador gráfico}, \textit{verificador dimensional}, \textit{portavoz}), andamiaje decreciente y debate plenario \citep{fullan2014rich,sawyer2008optimizing}.
  \item \textbf{Dimensiones del aprendizaje profundo:} \textit{Cognitiva} (contraste tabla--gráfica--ecuación y límites físicos), \textit{Interpersonal} (consenso argumentado y coevaluación entre pares) e \textit{Intrapersonal} (diario de metacognición tras cada sesión para cultivar mentalidad de crecimiento) \citep{schunk2012teorias}.
  \item \textbf{Evaluación formativa:} Recolección continua de predicción inicial, hoja de análisis por equipo, dictamen técnico y segunda versión revisada tras devolución descriptiva \citep{bostrom2023formative,eef2021feedback}.
\end{enumerate}
\end{tcolorbox}

La Figura~\ref{fig:sub06c-seq} detalla la secuencia didáctica de tres sesiones e indicadores formativos:

\begin{figure}[htbp]
\centering
\begin{tikzpicture}[font=\scriptsize, >=Latex, scale=0.92, transform shape]
  \node[draw=iiiepeNavy,fill=iiiepePale,thick,rounded corners=4pt,text width=2.8cm,minimum height=1.95cm,align=center] (f1) at (0,0) {\textbf{Sesión 1}\\Predicción inicial y análisis tabular\\(\textit{Saberes previos})};
  \node[draw=iiiepeBlue,fill=iiiepePale,thick,rounded corners=4pt,text width=2.8cm,minimum height=1.95cm,align=center] (f2) at (3.65,0) {\textbf{Sesión 2}\\Modelación gráfica y algebraica en tríos\\(\textit{Cognitivo + Social})};
  \node[draw=iiiepeAqua!80!black,fill=iiiepeMint,thick,rounded corners=4pt,text width=2.8cm,minimum height=1.95cm,align=center] (f3) at (7.3,0) {\textbf{Sesión 3}\\Debate de validez y retroalimentación\\(\textit{Evaluación formativa})};
  \node[draw=iiiepeOrange!90!black,fill=iiiepeWarm,thick,rounded corners=4pt,text width=2.8cm,minimum height=1.95cm,align=center] (f4) at (10.95,0) {\textbf{Cierre}\\Segunda versión y diario metacognitivo\\(\textit{Intrapersonal})};

  \draw[->,very thick,iiiepeNavy] (f1) -- (f2);
  \draw[->,very thick,iiiepeBlue] (f2) -- (f3);
  \draw[->,very thick,iiiepeAqua!80!black] (f3) -- (f4);
\end{tikzpicture}
\caption{Secuencia didáctica por realizar en tres sesiones de innovación matemática.}
\label{fig:sub06c-seq}
\end{figure}

{\small\setstretch{1.15}\renewcommand{\arraystretch}{1.3}
\begin{longtable}{|>{\raggedright\arraybackslash}p{0.22\linewidth}|>{\raggedright\arraybackslash}p{0.34\linewidth}|>{\raggedright\arraybackslash}p{0.38\linewidth}|}\hline
\rowcolor{iiiepePale} \textbf{Dimensión} & \textbf{Evidencia observable} & \textbf{Indicador formativo de avance} \\\hline
\endfirsthead
\hline
\rowcolor{iiiepePale} \textbf{Dimensión} & \textbf{Evidencia observable} & \textbf{Indicador formativo de avance} \\\hline
\endhead
Comprensión conceptual & Registros tabular--gráfico y ecuación justificada. & Distingue proporcional vs. afín e interpreta pendiente y ordenada en unidades correctas. \\\hline
Argumentación y uso & Dictamen escrito comparando rangos de validez. & Justifica decisiones señalando supuestos físicos y límites del modelo. \\\hline
Colaboración y mejora & Diario de corrección tras devolución descriptiva. & Integra observaciones de pares y corrige explícitamente errores previos \citep{eef2021feedback}. \\\hline
\end{longtable}}

\clearpage
\section{Conclusiones}
En conclusión, considero que esta propuesta operativa de innovación transforma los hallazgos de \citet{quiroga2022aprendizaje} y \citet{rivas2017innovaciones} en un protocolo realista, verificable mediante ciclos de investigación-acción y respetuoso de la heterogeneidad estudiantil. Sin embargo, dado que los datos de calibración presentados son didácticos simulados, además conviene contrastar su implementación real frente a producciones auténticas de grupo; es decir, por ello, por tanto sostengo su pertinencia para su futura aplicación e indagación empírica en el aula tecnológica y universitaria.

\clearpage
\begin{thebibliography}{99}
\bibitem[Boström y Palm(2023)]{bostrom2023formative} Boström, E., y Palm, T. (2023). Formative assessment practice in mathematics. \textit{Frontiers in Education}, 8, Artículo 1101192.
\bibitem[Collin y Quigley(2021)]{eef2021feedback} Collin, J., y Quigley, A. (2021). \textit{Teacher feedback to improve pupil learning}. Education Endowment Foundation.
\bibitem[CONALEP(2025)]{conalep2025mecatronica} Colegio Nacional de Educación Profesional Técnica. (2025). \textit{Perfil de egreso en Mecatrónica. Modelo Académico 2025}. CONALEP.
\bibitem[Fullan y Langworthy(2014)]{fullan2014rich} Fullan, M., y Langworthy, M. (2014). \textit{A rich seam: How new pedagogies find deep learning}. Pearson.
\bibitem[Quiroga Aguilar y Lara Órdenes(2022)]{quiroga2022aprendizaje} Quiroga Aguilar, L., y Lara Órdenes, E. (2022). El aprendizaje profundo como herramienta para cambio en la visión de aprendizaje. \textit{Revista Educación las Américas}, 12(1).
\bibitem[Rivas et~al.(2017)]{rivas2017innovaciones} Rivas, A., André, F., y Delgado, L. E. (Eds.). (2017). \textit{50 innovaciones educativas para escuelas}. CIPPEC / Santillana.
\bibitem[Sawyer(2008)]{sawyer2008optimizing} Sawyer, R. K. (2008). Optimizing learning: Implications of learning sciences research. OECD Publishing.
\bibitem[Schunk(2012)]{schunk2012teorias} Schunk, D. H. (2012). \textit{Teorías del aprendizaje} (6.ª ed.). Pearson.
\bibitem[SEP(2023)]{sep2023pensamiento} Secretaría de Educación Pública. (2023). \textit{Programa de estudios del recurso sociocognitivo Pensamiento Matemático I}. SEP.
\end{thebibliography}

\end{document}
